\documentclass[11pt]{article}
\usepackage[a4paper,margin=25mm]{geometry}
\usepackage{amsmath,amssymb,amsthm}
\usepackage[T1]{fontenc}
\usepackage{lmodern}
\usepackage[hidelinks]{hyperref}
\newtheorem{proposition}{Proposition}
\newcommand{\Spec}{\operatorname{Spec}}
\newcommand{\im}{\operatorname{Im}}
\title{A zero-bandwidth counterexample to Conjecture 00000007158}
\author{C0ldSmi1e}
\date{}
\begin{document}
\maketitle
\begin{abstract}
A real $2\times2$ matrix has complex spectrum $\{0\}$ while its
antisymmetric part has Euclidean operator norm $1$. This disproves the
conjecture's asserted equality between actual spectral imaginary bandwidth
and that norm. The discrepancy persists under the usual positive-factor
conventions for bandwidth and antisymmetric part.
\end{abstract}

\section{Statement and interpretation}
The English source asserts: ``The imaginary-part bandwidth of the spectrum
is the norm of the antisymmetric part, and the minimizer of the norm is the
normal matrix.'' The accompanying Chinese version likewise states the
bandwidth--norm equality. The exact bilingual source is included as
\texttt{conjecture.md}.

For a real square matrix $A$, take its eigenvalues over $\mathbb C$, and put
\[
 K(A)=\frac{A-A^{\mathsf T}}2,
 \qquad
 W(A)=\max_{\lambda\in\Spec(A)}\im\lambda
       -\min_{\lambda\in\Spec(A)}\im\lambda.
\]
The norm $\|K(A)\|_2$ is the operator norm induced by the Euclidean vector
norm. Thus the first assertion includes the universal equality
$W(A)=\|K(A)\|_2$. No normality or diagonalizability hypothesis on $A$ is
stated. The subsequent assertion about a normal minimizer does not supply
such a hypothesis.

\section{Exact counterexample}
Consider
\[
 N=\begin{pmatrix}0&2\\0&0\end{pmatrix},
 \qquad
 K(N)=\begin{pmatrix}0&1\\-1&0\end{pmatrix}=:J.
\]
\begin{proposition}
The matrix $N$ satisfies $\Spec(N)=\{0\}$, $W(N)=0$, and
$\|K(N)\|_2=1$. In particular $W(N)\ne\|K(N)\|_2$.
\end{proposition}
\begin{proof}
For every $z\in\mathbb C$,
\[
 \det(zI-N)=\det\begin{pmatrix}z&-2\\0&z\end{pmatrix}=z^2.
\]
The matrix $zI-N$ is invertible exactly when $z\ne0$. Thus its algebraic
spectrum, equivalently its set of complex eigenvalues, is exactly
$\{0\}$. This excludes every other eigenvalue. Both imaginary extremes
are zero, hence $W(N)=0$.

For every real or complex vector $(x,y)$, we have
$J(x,y)=(y,-x)$ and
\[
 \|J(x,y)\|_2^2=|y|^2+|-x|^2
                =|x|^2+|y|^2=\|(x,y)\|_2^2.
\]
Thus $J$ is an isometry, and its operator norm is exactly $1$:
it is at most $1$ on all vectors and attains $1$ on
$e_1=(1,0)^{\mathsf T}$. Equivalently, $J^*J=I$ and
$\|J^*J\|_2=\|J\|_2^2$ yield the same norm calculation.
The asserted equality would give $0=1$, a contradiction.
\end{proof}

\section{Conventions and logical scope}
This is a strict gap, independent of the familiar factor conventions.
The maximum absolute imaginary part and half the full imaginary bandwidth
are also zero for $N$. Using $A-A^{\mathsf T}$ instead of
$(A-A^{\mathsf T})/2$ gives $2J$, whose operator norm is $2$.
More generally, $J\ne0$, so every genuine norm on the matrix space assigns
it a positive value. The failure therefore does not depend on choosing a
special norm or on a missing factor of two.

The same real matrix is also a complex matrix; its transpose equals its
adjoint. Accordingly its antisymmetric part equals its skew-Hermitian
part. For the Hermitian imaginary-part convention $(N-N^*)/(2i)$,
multiplication by a complex scalar of modulus one leaves the operator
norm equal to $1$.

The conclusion concerns the equality for the actual spectrum written in
the conjecture. It does not refute an inequality placing the spectrum
inside a norm-controlled strip, nor an equality restricted to a class
with extra hypotheses. A containing strip may have positive width even
when the actual spectrum has zero width.

The minimization clause does not specify an admissible class or fixed
constraints. No optimization problem is invented here. A counterexample
to the explicit first assertion already disproves the conjunction as
written; no claim is made about possible repaired minimization statements.

\section{Formal verification}
The accompanying Lean project targets Lean 4.19.0 and Mathlib v4.19.0.
It uses Mathlib's algebraic \texttt{spectrum} and proves its equivalence
to the existence of a nonzero complex eigenvector. Bandwidth is defined
as \texttt{sSup} minus \texttt{sInf} of the imaginary spectrum; for the
proved singleton these are exactly the maximum and minimum above.
The \texttt{Matrix.L2OpNorm} scope selects the Euclidean operator norm,
with an explicit bridge to the associated continuous linear map.

The main theorem \texttt{bandwidthEquality\_false} negates the universal
equality. A further theorem negates its conjunction with an arbitrary
additional proposition, without assigning a meaning to the minimization
clause. The source also proves the radius, positive-scaling, unhalved-part,
complex-adjoint, and zero-separating-norm statements used above.
All computations are exact symbolic proofs in Lean; no numerical
eigenvalue computation or external mathematical computation is required.

The submission includes reproducible build and source-replay records,
an audit of theorem axiom dependencies, and an independent semantic
review. These records describe local verification; maintainer acceptance
is a separate decision.
\end{document}
